% Converted from Microsoft Word to LaTeX

% by Chikrii SoftLab Word2TeX converter (version 2.4)

% Copyright (C) 1999-2001 Kirill A. Chikrii, Anna V. Chikrii

% Copyright (C) 1999-2001 Chikrii SoftLab.

% All rights reserved.

% http://www.word2tex.com/

% mailto: info@word2tex.com, support@word2tex.com





\documentclass [12pt]{article}





\begin{document}

\begin{center}

\textbf{Об одном методе поиска глобального экстремума }

\end{center}



\begin{center}

\textbf{непрерывной функции на симплексе}

\end{center}











\begin{center}

Лившиц М.Ю., Сизиков А.П.

\end{center}











Рассматривается невыпуклая задача математического программирования, 

допустимой областью которой является симплекс. Для приближенного решения 

задачи предложен двухэтапный алгоритм. На первом этапе методом $\Psi - 

$преобразования определяется область глобального оптимума, на втором 

осуществляется локальная ``доводка'' решения. Метод $\Psi - $преобразования 

модифицирован с учетом специфики рассматриваемой задачи. Определение $\Psi - 

$функции осуществляется по результатам статистических испытаний, реализуемых 

с помощью генератора случайных чисел, равномерно распределенных в симплексе. 

Для уточнения решения применяется предложенный метод отражения правильных 

симплексов. Приведен пример применения разработанного алгоритма для решения 

задачи оптимизации компонентного состава углеводородной смеси.











\textbf{\textit{Ключевые слова}}: оптимизация, невыпуклые задачи, метод 

$\Psi - $преобразования, равномерное распределение в симплексе, 

многокомпонентные смеси. 











В широкой области технических приложений возникает необходимость решения 

нелинейной и, в общем случае, невыпуклой задачи: 





\begin{equation}

\label{eq1}

f(x) \to \max ,

\end{equation}









\begin{equation}

\label{eq2}

\left\{ {{\begin{array}{*{20}c}

 {\sum\limits_{j = 1}^m {x_j = 1,\;} } \\

 {x_j \ge 0,j = 1,2,...m,} \\

\end{array} }} \right.

\end{equation}







\noindent

где $f(x)_{ }$- некоторая непрерывная функция $x = (x_1 ,x_2 ,....x_m )$.



Для решения невыпуклых оптимизационных задач наиболее часто используют 

стохастический подход. Он включает широкий спектр методов - от простого 

случайного поиска до различных рандомизированных эвристик [1,2]. Все эти 

методы так или иначе используют глобальный и локальный поиск в различных 

реализациях и сочетаниях. 



В методе имитации отжига переход от глобального к локальному поиску 

осуществляется постепенно [3,4]. На начальном этапе осуществляется случайный 

поиск по всей допустимой области. Но затем, оставаясь случайным, он все 

более и более локализуется, становится избирательным. В методе 

кросс-энтропии [5-7] идея постепенной локализации поиска реализуется через 

механизм изменения распределения поисковой выборки с тем, чтобы 

перспективные варианты выпадали с большей вероятностью. Существуют методы, 

использующие так называемый ``мультистарт'' локальных поисков 

(детерминированных или рандомизированных) сразу из множества точек области 

определения [8,9]. Для решения некоторых трудно формализуемых задач 

эффективны генетические алгоритмы [10-12]. В целом же, эти алгоритмы 

обладают всеми достоинствами и недостатками рандомизированных эвристик [13].



Авторы предлагают для решения задачи (\ref{eq1}) - (\ref{eq2}) модификацию известного метода 

$\Psi - $преобразования [14, 15], сочетающего детерминированный и 

стохастический подходы. По сравнению с методами [1 - 12] предложенный метод 

менее чувствителен к размерности задачи. В его основе лежит разбиение, 

применяемое для построения интеграла Лебега. Для задачи $\max \left\{ 

{f(x)\vert x \in E} \right\}$, где$_{ }E$ - некоторое подмножество 

евклидова пространства, в рассмотрение вводится монотонно убывающая функция 

$\Psi (\zeta )$, которая$_{ }$определяется как мера Лебега области $E(\zeta 

) = \{x\vert x \in E,\,\,f(x) \ge \zeta \}$. Ноль функции $\Psi (\zeta )$ 

соответствует решению задачи. Определение функции $\Psi (\zeta )$ в 

аналитическом виде затруднительно, поэтому некоторое ее приближение строится 

по точкам, которые получают с помощью статистических испытаний. 



В поставленной задаче (\ref{eq1}) -- (\ref{eq2}) прежде чем использовать алгоритм $\Psi - 

$преобразования, необходимо при статистических испытаниях решить вопрос 

генерирования случайных точек, равномерно распределенных в области (\ref{eq2}). 



Правильный $(m - 1)$-мерный симплекс (\ref{eq2}), можно поместить в соответствующий 

$(m - 1)$-мерный единичный гиперкуб $R^{m - 1}$, получение равномерного 

распределения случайных чисел в котором -- задача тривиальная. Поскольку 

симплекс является подмножеством гиперкуба, точки, попадающие в симплекс, 

также равномерно распределены. Однако, для поставленной задачи (\ref{eq1}) - (\ref{eq2}) 

этот путь непродуктивен. Как известно, объем правильного $(m - 1)$-мерного 

симплекса со стороной единичной длины$_{ }$равен





\[

\frac{1}{(m - 1)!}\sqrt {\frac{m}{2^{m - 1}}} _{,}

\]







\noindent

и при росте $m_{ }$быстро стремится к нулю. Соответственно, стремится к 

нулю и вероятность того, что случайная точка, сгенерированная в гиперкубе, 

окажется в симплексе. Поэтому для получения случайных допустимых вариантов 

используем алгоритм [16, 17], который обеспечивает перенос равномерно 

распределенных случайных точек в симплекс (\ref{eq2}) путем аффинных и линейных 

преобразований, не влияющих на закон распределения.



Алгоритм 1:



1. Генерируются $(m - 1)_{ }$случайных чисел, равномерно распределенных на 

единичном отрезке. И эти числа нумеруются в порядке возрастания: $\xi _1 \le 

\xi _2 \le ... \le \xi _{m - 1} $. 



2. Формируется вектор ${\xi }' = \left( {{\xi }'_1 ,{\xi }'_2 ,...,{\xi 

}'_{m - 1} } \right),_{ }$где$_{ }{\xi }'_1 = \xi _1 $, ${\xi }'_2 = \xi 

_2 - \xi _1 $, \ldots , ${\xi }'_{m - 1} = \xi _{m - 1} - \xi _{m - 2} $. 

Этот вектор, как показано в работе [16], равномерно распределен на множестве 

$M = \{x \in R^{m - 1}\vert x_1 \ge 0,x_2 \ge 0,...,x_{m - 1} \ge 

0,\sum\limits_{i = 1}^{m - 1} {x_i \le 1} \}$.



3. Все угловые точки многогранника $M$, кроме нулевой, лежат на расстоянии 

$\sqrt 2 _{ }$друг от друга. Многогранник $M$ растягивается до правильного 

$\tilde {M}$. Для этого перенесем вершину, лежащую в начале координат так, 

чтобы расстояние от нее до всех остальных вершин было также $\sqrt 2 $. Это 

будет точка с координатами $(x_1^0 ,x_2^0 ,...,x_{m - 1}^0 )$, где $x_i^0 = 

r = (1 - \sqrt m ) / (m - 1),_{ }i = 1,2,...,m - 1.$ Получился 

многогранник, конгруэнтный $(m - 1)$-мерному симплексу в $m$-мерном 

пространстве. Обозначим через $\tilde {\xi }$точку множества $\tilde 

{M},_{ }$отображающую точку ${\xi }'_ _{ }$множества $M$. Точку $\tilde 

{\xi }_{ }$определяем в результате отражения ${\xi }'_ _{ }$относительно 

гиперплоскости $\{x \in R^{m - 1}\vert \sum\limits_{i = 1}^{m - 1} {x_i = 1} 

\}$. Для этого сначала находим проекцию ${\xi }'_{ }$на эту 

гиперплоскость: 

$\mathord{\buildrel{\lower3pt\hbox{$\scriptscriptstyle\smile$}}\over {\xi }} 

_i = {\xi }'_i + \Delta ,_{ }i = 1,2,...,m - 1,$ где $\Delta $определяем 

из условия $\sum\limits_{i = 1}^{m - 1} 

{\mathord{\buildrel{\lower3pt\hbox{$\scriptscriptstyle\smile$}}\over {\xi }} 

_i = 1} $ по формуле $\Delta = \frac{1}{m - 1}\left( {1 - \sum\limits_{i = 

1}^{m - 1} {{\xi }'_i } } \right)$.$_{ }$Отражение находим с учетом того, 

что коэффициент растяжения равен $\sqrt m $: $\tilde {\xi } = 

\mathord{\buildrel{\lower3pt\hbox{$\scriptscriptstyle\smile$}}\over {\xi }} 

+ \sqrt т ({\xi }' - 

\mathord{\buildrel{\lower3pt\hbox{$\scriptscriptstyle\smile$}}\over {\xi }} 

)$. Это аффинное преобразование не оказывает влияние на закон распределения, 

поэтому точка $\tilde {\xi }$имеет равномерное распределение в $\tilde {M}$.



4. Точка $\tilde {\xi }_{ }$переносится в $(m - 1)$- мерный симплекс $m$- 

мерного пространства переменных задачи (\ref{eq1}) - (\ref{eq2}). Для этого представляем ее 

как выпуклую линейную комбинацию угловых точек $\tilde {M}_{ }m$- 

мерного пространства. Коэффициенты этой комбинации $\alpha _1 ,\alpha _2 

,...,\alpha _m _{ }$есть не что иное, как новые координаты точки $\alpha 

$. Они определяются следующим образом:





\[

\left( {{\begin{array}{*{20}c}

 {\alpha _1 } \hfill \\

 {\alpha _2 } \hfill \\

 {...} \hfill \\

 {\alpha _{m - 1} } \hfill \\

 {\alpha _m } \hfill \\

\end{array} }} \right) \quad _{ = }

\quad

\left( {{\begin{array}{*{20}c}

 1 \hfill & 0 \hfill & {...} \hfill & 0 \hfill & r \hfill \\

 0 \hfill & 1 \hfill & {...} \hfill & 0 \hfill & r \hfill \\

 {...} \hfill & {...} \hfill & {...} \hfill & {...} \hfill & {...} \hfill \\

 0 \hfill & 0 \hfill & {...} \hfill & 1 \hfill & r \hfill \\

 1 \hfill & 1 \hfill & {...} \hfill & 1 \hfill & 1 \hfill \\

\end{array} }} \right)^{ - 1}

\left( {{\begin{array}{*{20}c}

 {\tilde {\xi }_1 } \hfill \\

 {\tilde {\xi }_2 } \hfill \\

 {...} \hfill \\

 {\tilde {\xi }_{m - 1} } \hfill \\

 1 \hfill \\

\end{array} }} \right)_{.}

\]







Теперь вернемся к методу $\Psi - $преобразования. 



Алгоритм 2:



1. Выполняется серия случайных проб, число которых $N$должно быть 

достаточным для получения надежных статистических оценок. На каждой пробе $k 

= 1,2,...N_{ }$с помощью алгоритма 1 генерируется допустимое значение 

аргумента $x^{(k)}_{ }$и находится $f(x^{(k)})$. По результатам всех проб 

определяются среднее арифметическое $f_c $ и максимальное $f_m $ наблюдаемых 

значений целевой функции: 





\[

f_c = \frac{1}{N}\sum\limits_{k = 1}^N {f(x^{(k)})} ,

\quad

f_m = \max \left\{ {f(x^{(\ref{eq1})}),f(x^{(\ref{eq2})}),...,f(x^{(N)})} \right\}.

\]







Отрезок $[f_c ,f_m ]$ разбивается на $L_{ }$(порядка десяти) уровней 

$\zeta = (\zeta _1 ,\zeta _2 ,...,\zeta _L )$ где $\zeta _l = f_c + (l - 

1)(f_m - f_c )L^{ - 1}$. 



2. Выполняется вторая серия случайных проб. Для каждого уровня $l = 

1,2,...,L_{ }$подсчитывается$_{ }n_l _{ }$- число проб, в которых 

выполнилось условие $f(x^{(k)}) \ge \zeta _l $, и средняя координата: 





\[

\bar {x}_{il} = \frac{1}{n_l }\sum\limits_{k = 1}^N {\delta _l } x_i^{(k)} 

,i = 1,2,...,m,\delta _l = \left\{ {{\begin{array}{*{20}c}

 {1,f(x^{(k)}) \ge \zeta _l ,} \hfill \\

 {0,f(x^{(k)}) < \zeta _l .} \hfill \\

\end{array} }} \right.

\]







По результатам формируются последовательности: 





\begin{equation}

\label{eq3}

\psi = \{\psi _l \}_{l = 1}^L ,\mbox{г}\mbox{д}\mbox{е}

\quad

\psi _l = n_l N^{ - 1},

\end{equation}









\begin{equation}

\label{eq4}

\tilde {x}_i = \{\bar {x}_{il}^ \}_{l = 1}^L ,

\quad

i = 1,2,...,m.

\end{equation}







3. Определяются квадратичные тренды. Пусть 





\[

\Lambda = \left( {{\begin{array}{*{20}c}

 1 \hfill & 1 \hfill & 1 \hfill \\

 1 \hfill & 2 \hfill & 4 \hfill \\

 {...} \hfill & {...} \hfill & {...} \hfill \\

 1 \hfill & L \hfill & {L^2} \hfill \\

\end{array} }} \right) \quad ,

\]







\noindent

тогда параметры $a = (a_0 ,a_1 ,a_2 )^T_{ }$тренда $\psi (l) = a_0 + a_1 l 

+ a_2 l^2$ рассчитываются по формуле $a = (\Lambda ^T\Lambda )^{ - 1}\Lambda 

^T\psi ^T$. Аналогично, параметры $b_i = (b_{i0} ,b_{i1} ,b_{i2} )^T$ тренда 

$\tilde {x}_i (l) = b_{i0} + b_{i1} l + b_{i2} l^2_{ }$- по формуле $b_i = 

(\Lambda ^T\Lambda )^{ - 1}\Lambda ^T\tilde {x}_i^T $.



4. Определяется наименьший положительный корень $l^\ast _{ }$уравнения 

$\psi (l) = 0$. Его существование обусловлено тем, что $\psi (l)_{ }$на 

отрезке $[1,L]_{ }$представляет собой тренд последовательности (\ref{eq3}), 

монотонно убывающей до значения, близкого к нулю. Подстановкой $l^\ast _{ 

}$в тренды $\tilde {x}(l)$ находится приближенное решение задачи: $x_i^\ast 

= \tilde {x}_i (l^\ast )$, $i = 1,2,...,m$.



Назначение описанного алгоритма 2 -- войти в область глобального экстремума. 

Предположим, что это удалось. Перейдем к локальному поиску самого 

экстремума. Методы, предполагающие использование производных первого или 

второго порядков (метод сопряженных градиентов, метод Ньютона) здесь не 

применимы по причине возможной негладкости целевой функции $f(x)$. В этом 

случае необходимо использовать метод, для реализации которого достаточно 

значений самой целевой функции $f(x)$. Например, последовательное отражение 

симплексов [18]. 



Классический метод последовательного отражения симплексов -- это метод 

поиска экстремума в $m$-мерном евклидовом пространстве с помощью $m$-мерных 

симплексов [18]. В задаче (\ref{eq1}) - (\ref{eq2}) речь идет о поиске экстремума в $(m - 

1)$-мерном симплексе с помощью $(m - 1)$-мерных же симплексов.



Построим в $m$-мерном евклидовом пространстве правильный $m$-мерный 

симплекс, одна из вершин которого (назовем ее нулевой) лежит в начале 

координат. Пусть длина ребра равна $\Delta x$, тогда координаты вершин 

соответствуют столбцам матрицы:





\[

\left( {{\begin{array}{*{20}c}

 {x_1^{(0)} } \hfill & {x_1^{(\ref{eq1})} } \hfill & {x_1^{(\ref{eq2})} } \hfill & {...} 

\hfill & {x_1^{(m)} } \hfill \\

 {x_2^{(0)} } \hfill & {x_2^{(\ref{eq1})} } \hfill & {x_2^{(\ref{eq2})} } \hfill & {...} 

\hfill & {x_2^{(m)} } \hfill \\

 {...} \hfill & {...} \hfill & {...} \hfill & {...} \hfill & {...} \hfill \\

 {x_m^{(0)} } \hfill & {x_m^{(\ref{eq1})} } \hfill & {x_m^{(\ref{eq2})} } \hfill & {...} 

\hfill & {x_m^{(m)} } \hfill \\

\end{array} }} \right) \quad = 

\quad

\left( {{\begin{array}{*{20}c}

 0 \hfill & 

\mathord{\buildrel{\lower3pt\hbox{$\scriptscriptstyle\frown$}}\over {q}} 

\hfill & q \hfill & {...} \hfill & q \hfill \\

 0 \hfill & q \hfill & 

\mathord{\buildrel{\lower3pt\hbox{$\scriptscriptstyle\frown$}}\over {q}} 

\hfill & {...} \hfill & q \hfill \\

 {...} \hfill & {...} \hfill & {...} \hfill & {...} \hfill & {...} \hfill \\

 0 \hfill & q \hfill & q \hfill & {...} \hfill & 

\mathord{\buildrel{\lower3pt\hbox{$\scriptscriptstyle\frown$}}\over {q}} 

\hfill \\

\end{array} }} \right),

\]







\noindent

где $\mathord{\buildrel{\lower3pt\hbox{$\scriptscriptstyle\frown$}}\over 

{q}} = \frac{\Delta x}{m\sqrt 2 }\left( {\sqrt {m + 1} + m - 1} \right)$, $q 

= \frac{\Delta x}{m\sqrt 2 }\left( {\sqrt {m + 1} - 1} \right)$.



Перенесем этот симплекс из начала координат так, чтобы все его вершины, 

кроме $x^{(0)}_{ }$расположились на гиперплоскости $x_1 + x_2 + ... + x_m 

= 1$ и чтобы $(m - 1)$-мерный симплекс с вершинами 

$x^{(\ref{eq1})},x^{(\ref{eq2})},...,x^{(m)}$ оказался в области начала поиска. 



При параллельном переносе нулевой вершины в точку $\tilde {x}^{(0)}$ 

координаты остальных вершин становятся равными $\tilde {x}^{(k)} = \tilde 

{x}^{(0)} + x^{(k)}$, $k = 1,2,...,m$. Найдем $\tilde {x}^{(0)}_{ }$из 

условия $\tilde {x}_1^{(k)} + \tilde {x}_2^{(k)} + ... + \tilde {x}_m^{(k)} 

= 1$. Достаточно взять какую-нибудь одну вершину, например, первую. 





\[

\sum\limits_{i = 1}^m {\tilde {x}_i^{(\ref{eq1})} = } 

\sum\limits_{i = 1}^m {\tilde {x}_i^{(0)} } 

 + \frac{\Delta x}{m\sqrt 2 }\left( {\sqrt {m + 1} + m - 1} \right)

 + \sum\limits_{i = 2}^m {\frac{\Delta x}{m\sqrt 2 }\left( {\sqrt {m + 1} - 

1} \right)} \quad = 

\]







\begin{center}

= $\sum\limits_{i = 1}^m {\tilde {x}_i^{(0)} }  + \frac{\Delta x}{m\sqrt 2 

}\left( {\sqrt {m + 1} + m - 1 + \sum\limits_{i = 1}^{m - 1} {\left( {\sqrt 

{m + 1} - 1} \right)} } \right)_{ }$=

\end{center}



\begin{center}

= $\sum\limits_{i = 1}^m {\tilde {x}_i^{(0)} }  + \frac{\Delta x}{m\sqrt 2 

}m\sqrt {m + 1} _{ }$= 1.

\end{center}



Отсюда получаем: 





\[

\sum\limits_{i = 1}^m {\tilde {x}_i^{(0)} } = 

\quad

1 - \frac{\Delta x}{\sqrt 2 }\sqrt {m + 1} .

\]







Определим начальное расположение $m$ -- мерного симплекса, привязав его к точке 

$x^\ast $, полученной методом $\Psi - $преобразования. В этом случае 

координаты нулевой вершины симплекса определятся следующим образом:





\[

\tilde {x}_i^{(0)} = 

\quad

x_i^\ast \left( {1 - \frac{\Delta x}{\sqrt 2 }\sqrt {m + 1} } \right),

\quad

i = 1,2,...,m.

\]







Координаты остальных вершин, которые и образуют начальный ``рабочий'' 

симплекс, представлены столбцами матрицы:





\[

\left( {{\begin{array}{*{20}c}

 {\tilde {x}_1^{(\ref{eq1})} } \hfill & {\tilde {x}_1^{(\ref{eq2})} } \hfill & {...} \hfill 

& {\tilde {x}_1^{(m)} } \hfill \\

 {\tilde {x}_2^{(\ref{eq1})} } \hfill & {\tilde {x}_2^{(\ref{eq2})} } \hfill & {...} \hfill 

& {\tilde {x}_2^{(m)} } \hfill \\

 {...} \hfill & {...} \hfill & {...} \hfill & {...} \hfill \\

 {\tilde {x}_m^{(\ref{eq1})} } \hfill & {\tilde {x}_m^{(\ref{eq2})} } \hfill & {...} \hfill 

& {\tilde {x}_m^{(m)} } \hfill \\

\end{array} }} \right),

\]







\noindent

где





\[

\tilde {x}_i^{(k)} = 

\quad

x_i^\ast \left( {1 - \frac{\Delta x}{\sqrt 2 }\sqrt {m + 1} } \right) + 

\left\{ {{\begin{array}{*{20}c}

 {\frac{\Delta x}{m\sqrt 2 }\left( {\sqrt {m + 1} + m - 1} \right),i = k,} 

\\

 {\frac{\Delta x}{m\sqrt 2 }\left( {\sqrt {m + 1} - 1} \right),i \ne k.} \\

\end{array} }} \right.

\]







Процесс поиска состоит в реализации последовательности шагов, на каждом из 

которых происходит построение нового рабочего симплекса путем отражения 

вершины с минимальным значением целевой функции $f(x)_{ }$от 

противоположной грани исходного симплекса. Пусть $v \in [1,m]_{ }$- номер 

такой вершины, тогда отраженная вершина определяется по формуле:





\[

\mathord{\buildrel{\lower3pt\hbox{$\scriptscriptstyle\smile$}}\over {x}} 

^{(v)} = \frac{2}{m - 1}\left( {\sum\limits_{k = 1}^m {x_^{(k)} - x^{(v)}} } 

\right) - x^{(v)}.

\]







Покажем, что 

$\mathord{\buildrel{\lower3pt\hbox{$\scriptscriptstyle\smile$}}\over {x}} 

^{(v)}_{ }$лежит на гиперплоскости $x_1 + x_2 + ... + x_m = 1$. 

Действительно,





\[

\sum\limits_{i = 1}^m 

{\mathord{\buildrel{\lower3pt\hbox{$\scriptscriptstyle\smile$}}\over {x}} 

_i^{(v)} } = \frac{2}{m - 1}\sum\limits_{i = 1}^m {\left( {\sum\limits_{k = 

1}^m {x_i^{(k)} - x_i^{(v)} } } \right)} - \sum\limits_{i = 1}^m {x_{i = 

1}^{(v)} } \quad = 

\]









\[

\frac{2}{m - 1}\left( {\sum\limits_{k = 1}^m {\sum\limits_{i = 1}^m 

{x_i^{(k)} - \sum\limits_{i = 1}^m {x_i^{(v)} } } } } \right) - 

\sum\limits_{i = 1}^m {x_i^{(v)} } = 

\frac{2}{m - 1}\left( {m - 1} \right) - 1 = 1.

\]







Поскольку поиск проводится в ограниченном пространстве, может оказаться, что 

$\mathord{\buildrel{\lower3pt\hbox{$\scriptscriptstyle\smile$}}\over {x}} 

^{(v)}$ выходит за допустимую область. Предположим для простоты, что на 

текущем шаге алгоритма - $f\left( {x^{(\ref{eq1})}} \right) < f\left( {x^{(\ref{eq2})}} 

\right)... < f\left( {x^{(m)}} \right)$. Тогда нужно последовательно 

перебирать вершины до тех пор, пока не получим допустимого отражения. 

Отражение лучшей точки будет признаком окончания текущей серии. Следующая 

серия начинается с уменьшения длины ребра $\Delta x_{ }$и переноса нового 

``рабочего'' симплекса в область лучшей точки предыдущей серии. 



Описанная методика может быть использована для поиска оптимального 

компонентного состава различных смесей, в частности, для расчета рецептур 

смешения товарных нефтепродуктов. Некоторые показатели качества смесей 

рассчитываются по специальным, иногда довольно сложным эмпирическим формулам 

[19, 20]. 



Пусть $Q$ - множество индексов контролируемых показателей некоторой смеси, и 

пусть $\varphi _q (x)_{ }$- зависимость (линейная или нелинейная) $q$-го 

показателя смеси от компонентов смешения, взятых в пропорциях $x = (x_1 ,x_2 

,....x_m )$. Тогда задачу смешения можно сформулировать так:





\[

C = \sum\limits_{j = 1}^m {c_j x_j } \to \min ,

\]









\[

\left\{ {{\begin{array}{*{20}c}

 {\phi _q (x) \in [p_q^ - ,p_q^ + ],q \in Q,} \\

 {\sum\limits_{j = 1}^m {x_j = 1,} } \\

 {x_j \ge 0,j = 1,2,...m,} \\

\end{array} }} \right.

\]







\noindent

где $c_j _{ }$- цена $j$-го компонента,$_{ }p_q^ - ,p_q^ + $ - граничные 

значения для $q$-го показателя.



Требования по качеству смеси удобно учесть в виде штрафной составляющей 

целевой функции: 





\begin{equation}

\label{eq5}

f(x) = \sum\limits_{j = 1}^m {c_j x_j + \mu \sum\limits_{q \in Q} {\varphi 

_q^2 (x)} } \to \min ,

\end{equation}









\begin{equation}

\label{eq6}

\left\{ {{\begin{array}{*{20}c}

 {\sum\limits_{j = 1}^m {x_j = 1,} } \\

 {x_j \ge 0,j = 1,2,...m,} \\

\end{array} }} \right.

\end{equation}







\noindent

где $\mu $- весовой коэффициент штрафной составляющей;$_{ }\varphi _q 

(x)$- процентное отклонение расчетного значения $q$-го показателя от границы 

заданного интервала, определяемое по формуле: 





\[

\varphi _q (x) = \left\{ {{\begin{array}{*{20}c}

 {10^2\left( {p_q^ - - \phi _q (x)} \right) / p_q^ - ,p_q^ - - \phi _q (x) > 

0,} \\

 {10^2\left( {\phi _q (x) - p_q^ + } \right) / p_q^ + ,p_q^ + - \phi _q (x) 

< 0,} \\

 {0,\phi _q (x) \in [p_q^ - ,p_q^ + ].} \\

\end{array} }} \right.

\]















Для примера решена задача составления трехкомпонентной смеси со следующими 

условиями. Стоимость $C = 1,5x_1 + 1,3x_2 + 1,0x_3 $. Требования к 

показателям (октановому числу, плотности, содержанию серы) : 



$\left\{ {{\begin{array}{*{20}c}

 {\phi _1 (x) = 90,2x_1 + 88,5x_2 + 73,6x_3 - 14,3x_1 x_2 - 9,8x_1 x_3 + 

28,4x_1 x_2 x_3 \ge 85,} \\

 {\phi _2 (x) = 0,82x_1 + 0,59x_2 + 0,67x_3 \le 0,68,} \\

 {\phi _3 (x) = 0,001x_1 + 0,003x_2 + 0,002x_3 \le 0,002.} \\

\end{array} }} \right._{(7)



При составлении функции $f(x)$ в (\ref{eq5}) используется весовой коэффициент $\mu = 

10$. По результатам предварительного статистического исследования области 

изменения $f(x)_{ }$сформировано 10 уровней с шагом $\Delta \zeta = - 70$, 

начиная с $\zeta _1 = 700$. 



Для каждого из них выполнено 1000 статистических испытаний. Получены 

последовательности (\ref{eq3}), (\ref{eq4}) и построены тренды:



\noindent

a) $\psi (l) = 0,6916 - 0,02906l - 0,00374l^2$,



\noindent

b) $\tilde {x}_1 (l) = 0,33041 - 0,00145l + 0,000098l^2$,



\noindent

c)$\tilde {x}_2 (l) = 0,327667 + 0,013837l - 0,00072l^2$,



\noindent

d) $\tilde {x}_3 (l) = 0,341923 - 0,012392l + 0,000626l^2$.





\begin{table}[htbp]

\begin{tabular}

{|p{225pt}|p{225pt}|}

\hline

 & 

  \\

\hline

 & 

  \\

\hline

\multicolumn{2}{|p{451pt}|}{$_{\mbox{Т}\mbox{р}\mbox{е}\mbox{н}\mbox{д}\mbox{ы} \mbox{р}\mbox{е}\mbox{з}\mbox{у}\mbox{л}\mbox{ь}\mbox{т}\mbox{а}\mbox{т}\mbox{о}\mbox{в} \mbox{с}\mbox{т}\mbox{а}\mbox{т}\mbox{и}\mbox{с}\mbox{т}\mbox{и}\mbox{ч}\mbox{е}\mbox{с}\mbox{к}\mbox{и}\mbox{х} \mbox{и}\mbox{с}\mbox{п}\mbox{ы}\mbox{т}\mbox{а}\mbox{н}\mbox{и}\mbox{й} }$}  \\

\hline

\end{tabular}

\label{tab1}

\end{table}















Из уравнения $\psi (l) = a_0 + a_1 l + a_2 l^2 = 0$ получено $l^\ast = 

10,2554.$ Подстановкой $l^\ast $ в тренды $\tilde {x}_1 (l)$, $\tilde {x}_2 

(l)$, $\tilde {x}_3 (l)_{ }$. найдено приближенное решение: $x_1^\ast = 

0,3259, \quad x_2^\ast = 0,3934,_{ }x_3^\ast = 0,2807.$ Значение целевой 

функции (\ref{eq5}):$f(x^\ast ) = 115,3.$



Уточнение решения осуществляется по описанному выше модифицированному методу 

отражения правильных симплексов. Проведено 20 серий с уменьшением длины 

ребра симплекса по формуле $\Delta x_n = 0,1n^{ - 1},$ где $n_{ }$- номер 

серии. Результаты первого приближения $x_^\ast _{ }$и уточненного 

$\mathord{\buildrel{\lower3pt\hbox{$\scriptscriptstyle\smile$}}\over {x}} 

_{ }$решения представлены в таблице









\begin{table}[htbp]

\begin{tabular}

{|p{35pt}|p{57pt}|p{57pt}|p{57pt}|p{57pt}|p{57pt}|p{57pt}|p{57pt}|}

\hline

& 

$x_1 $& 

$x_2 $& 

$x_3 $& 

$\phi _1 $& 

$\phi _2 $& 

$\phi _3 $& 

$f(x)$ \\

\hline

$x_^\ast $& 

0,3259& 

0,3934& 

0,2807& 

83,16& 

0,687& 

0,00209& 

115,3 \\

\hline

$\mathord{\buildrel{\lower3pt\hbox{$\scriptscriptstyle\smile$}}\over {x}} $& 

0,3702& 

0,3990& 

0,2308& 

83,71& 

0,693& 

0,00206& 

87,6 \\

\hline

\end{tabular}

\label{tab2}

\end{table}











В результате реализации алгоритма выявлена несовместность системы 

ограничений (7) и получено решение по критерию (\ref{eq5}), основная составляющая 

которого - взвешенная сумма квадратов процентных отклонений показателей от 

требуемых значений.











БИБЛИОГРАФИЧЕСКИЙ СПИСОК











1 Zhigljavsky A., Zilinskas A. Stochastic Global Optimization. Springer 

Science-Business Media, 2008. - 262 p. 



2. Пантелеев А.В. Метаэвристические алгоритмы поиска глобального экстремума. 

- М.: Издательство МАИ-ПРИНТ, 2009. -- 159 с.



3. Джонс М. Т. Программирование искусственного интеллекта в приложениях. - 

М.: ДМК Пресс, 2004. - 312 с. 



4. Савин А. Н., Тимофеева Н. Е. Применение алгоритма оптимизации методом 

имитации отжига на системах параллельных и распределённых вычислений. 

Известия Саратовского университета. Новая серия. Серия Математика. Механика. 

Информатика. -- Саратов: Изд-во Сарат. ун-та, 2012. -- Том 12. Вып. 1. -- С 

110-116. 



5. Z. I. Botev and D. P. Kroese. Global likelihood optimization via the 

crossentropy method with an application to mixture models. In Proceedings of 

the 36th Winter Simulation Conference, pages 529--535, Washington, D.C., 

2004.



6. D. Ernst, M. Glavic, G.-B. Stan, S. Mannor, L. Wehenkel, and Gif sur 

Yvette Supelec. The cross-entropy method for power system combinatorial 

optimization problems. In Power Tech, 2007 IEEE Lausanne, pages 1290 -- 

1295, Lausanne, 2007. 



7. G. E. Evans, J. M. Keith, and D. P. Kroese. Parallel cross-entropy 

optimization. In Proceedings of the 2007 Winter Simulation Conference, pages 

2196--2202, Washington, D.C., 2007.



8. Жиглявский А.А., Жилинская А.Г. Методы поиска глобального экстремума. - 

М.: Наука, 1991. - 248 с



9. Феоктистов А. Г., Горский С. А. Реализация метода мультистарта в пакете 

Градиент// Вестник НГУ Серия: Информационные технологии. - 2007. - Том 05, 

Выпуск № 2. - С. 78-82. 



10. Cohoon J.P., Karro J., Lienig J. Evolutionary Algorithms for the 

Physical Design of VLSI Circuits. Advances in Evolutionary Computing: Theory 

and Applications, Ghosh, A., Tsutsui, S. (eds.) Springer Verlag, London, 

2003. -- P. 683-712. 



11. Гладков Л.А., Курейчик В.В., Курейчик В.М. Генетические алгоритмы/ Под 

ред. В.М. Курейчика. - 2-е изд., испр. и доп. - М.: ФИЗМАТЛИТ, 2006. - 320 

с. 



12. Гладков Л.А., Гладкова Н.В. Особенности использования нечетких 

генетических алгоритмов для решения задач оптимизации и управления // 

Известия ЮФУ. Технические науки. - 2009. № 4 (93). -- C. 130-136.











13. Steven S. Skiena. The Algorithm Design Manual. Second Edition. Springer, 

2008. -- 730 p.



14. В.К. Чичинадзе. Решение невыпуклых нелинейных задач оптимизации. Метод 

$\Psi - $преобразования. Москва ``Наука''. Главная редакция 

физико-математической литературы. 1983. -- 256 с.



15. Ахмадиев Ф.Г., Гильфанов Р.М. Математическое моделирование и оптимизация 

``состав-свойство'' многокомпонентных смесей. Математическое моделирование, 

численные методы и комплексы программ (в строительстве) Известия КГАСУ, 

2012, № 2 (20)



16. Новоселов А.А. Равномерное распределение на стандартном симплексе. 

[Электронный ресурс] // Электронная библиотека Попечительского совета 

механико-математического факультета Московского государственного 

университета [Офиц. сайт]. URL: http://lib.mexmat.ru/books/31728



17. Ю.В. Русин Алгоритмы статистического моделирования вероятностных 

распределений: текст лекций/ Яросл. гос. ун-т. Ярославль: ЯрГУ, 2006.



18. Зедгинидзе И.Г. Планирование эксперимента для исследования 

многокомпонентных систем. Москва, ``Наука'', 1976. -- 390 с. 



19. Чинакал В.О. Оптимизация рецептуры светлых нефтепродуктов. -- В кн.: 

Оптимизация, исследование операций, бионика. М.: ``Наука'', 1973, С. 198 -- 

205. 



20. Никитин В.А., Мусаев А.А.. Оптимизация компаундирования углеводородных 

смесей // Труды СПИИРАН. 2007. Вып. 4. C. 327-336.











\end{document}

